\documentclass[12pt, a4paper]{ctexart}
\usepackage{amsmath, amssymb, amsthm}
\usepackage{geometry}
\usepackage{fancyhdr}
\usepackage{enumitem}
\usepackage{tikz}
\usepackage{pgfplots}
\usepackage{array}
\usepackage{multirow}
\usepackage{booktabs}

\geometry{left=2cm, right=2cm, top=2.5cm, bottom=2.5cm}

\setCJKmainfont{Hiragino Sans GB}
\setCJKsansfont{Heiti SC}

\pagestyle{fancy}
\fancyhf{}
\fancyhead[L]{一元函数的导数及其应用 测试卷}
\fancyhead[R]{数学（选择性必修第二册）}
\fancyfoot[C]{第 \thepage\ 页}
\renewcommand{\headrulewidth}{0.4pt}

\setlength{\parindent}{0pt}

\newcommand{\score}[1]{\hfill\textbf{（#1分）}}

\begin{document}

\begin{center}
{\Large\bfseries 一元函数的导数及其应用 测试卷}\\[0.5em]
{\normalsize 时间：120分钟 \quad 满分：150分}
\end{center}

\vspace{1em}

\noindent\textbf{一、单项选择题（本大题共8小题，每小题5分，共40分。每小题只有一个选项符合题目要求。）}

\vspace{0.5em}

\begin{enumerate}
\item 函数 $f(x)=x^2$ 在 $x=1$ 处的导数为 \score{5}
\begin{enumerate}[label=\Alph*.]
\item $1$
\item $2$
\item $x$
\item $2x$
\end{enumerate}

\item 函数 $f(x)=3x^3-2x+1$ 的导数 $f'(x)=$ \score{5}
\begin{enumerate}[label=\Alph*.]
\item $9x^2-2$
\item $9x^2-2x$
\item $3x^2-2$
\item $9x^2+1$
\end{enumerate}

\item 函数 $f(x)=\sin x$ 的导数 $f'(x)=$ \score{5}
\begin{enumerate}[label=\Alph*.]
\item $\cos x$
\item $-\cos x$
\item $\sin x$
\item $-\sin x$
\end{enumerate}

\item 曲线 $y=x^3-3x+1$ 在点 $(1,-1)$ 处的切线斜率为 \score{5}
\begin{enumerate}[label=\Alph*.]
\item $0$
\item $1$
\item $2$
\item $3$
\end{enumerate}

\item 函数 $f(x)=x^3-3x^2+1$ 的单调递增区间为 \score{5}
\begin{enumerate}[label=\Alph*.]
\item $(-\infty,0)$
\item $(2,+\infty)$
\item $(-\infty,0)\cup(2,+\infty)$
\item $(0,2)$
\end{enumerate}

\item 函数 $f(x)=x^3-3x$ 在 $x=-1$ 处取得 \score{5}
\begin{enumerate}[label=\Alph*.]
\item 极大值 $2$
\item 极小值 $2$
\item 极大值 $-2$
\item 极小值 $-2$
\end{enumerate}

\item 函数 $f(x)=e^x-x$ 的最小值为 \score{5}
\begin{enumerate}[label=\Alph*.]
\item $0$
\item $1$
\item $e$
\item $e-1$
\end{enumerate}

\item 函数 $f(x)=\ln x-\dfrac{1}{x}$ 在区间 $[1,e]$ 上的最大值为 \score{5}
\begin{enumerate}[label=\Alph*.]
\item $0$
\item $1-\dfrac{1}{e}$
\item $1$
\item $\ln e-\dfrac{1}{e}$
\end{enumerate}
\end{enumerate}

\vspace{1em}

\noindent\textbf{二、多项选择题（本大题共4小题，每小题5分，共20分。每小题有多个选项符合题目要求，全部选对得5分，部分选对得2分，有选错的得0分。）}

\vspace{0.5em}

\begin{enumerate}[resume]
\item 关于导数，下列结论正确的是 \score{5}
\begin{enumerate}[label=\Alph*.]
\item $(C)'=0$（$C$ 为常数）
\item $(x^n)'=nx^{n-1}$
\item $(\sin x)'=\cos x$
\item $(\cos x)'=\sin x$
\end{enumerate}

\item 函数 $f(x)=x^3-12x$，则下列结论正确的是 \score{5}
\begin{enumerate}[label=\Alph*.]
\item $f(x)$ 在 $(-\infty,-2)$ 上单调递增
\item $f(x)$ 在 $(-2,2)$ 上单调递减
\item $f(x)$ 在 $(2,+\infty)$ 上单调递增
\item $f(x)$ 的极大值为 $16$
\end{enumerate}

\item 关于函数 $f(x)=x-\ln x$（$x>0$），下列结论正确的是 \score{5}
\begin{enumerate}[label=\Alph*.]
\item $f(x)$ 在 $(0,1)$ 上单调递减
\item $f(x)$ 在 $(1,+\infty)$ 上单调递增
\item $f(x)$ 的最小值为 $1$
\item $f(x)\geq 1$ 恒成立
\end{enumerate}

\item 已知函数 $f(x)=\dfrac{\ln x}{x}$（$x>0$），则下列结论正确的是 \score{5}
\begin{enumerate}[label=\Alph*.]
\item $f(x)$ 在 $(0,e)$ 上单调递增
\item $f(x)$ 在 $(e,+\infty)$ 上单调递减
\item $f(x)$ 的最大值为 $\dfrac{1}{e}$
\item $f(e)=f(e^3)$
\end{enumerate}
\end{enumerate}

\vspace{1em}

\noindent\textbf{三、填空题（本大题共4小题，每小题5分，共20分。）}

\vspace{0.5em}

\begin{enumerate}[resume]
\item 函数 $f(x)=(2x+1)^3$ 的导数 $f'(x)=$ \underline{\hspace{3cm}}。\score{5}

\item 函数 $f(x)=xe^x$ 的导数 $f'(x)=$ \underline{\hspace{3cm}}。\score{5}

\item 函数 $f(x)=x^3-3x^2+2$ 在区间 $[-1,3]$ 上的最大值为 \underline{\hspace{2cm}}，最小值为 \underline{\hspace{2cm}}。\score{5}

\item 曲线 $y=x^3+x$ 在点 $(1,2)$ 处的切线方程为 \underline{\hspace{5cm}}。\score{5}
\end{enumerate}

\vspace{1em}

\noindent\textbf{四、解答题（本大题共6小题，共70分。解答应写出文字说明、证明过程或演算步骤。）}

\vspace{0.5em}

\begin{enumerate}[resume]
\item （10分）已知函数 $f(x)=x^3-3x^2-9x+5$。
\begin{enumerate}[label=(\arabic*)]
\item 求 $f'(x)$；
\item 求 $f(x)$ 的单调区间和极值。
\end{enumerate}

\item （10分）已知函数 $f(x)=x-\dfrac{1}{x}-2\ln x$（$x>0$）。
\begin{enumerate}[label=(\arabic*)]
\item 求 $f'(x)$；
\item 求 $f(x)$ 的单调区间；
\item 证明：当 $x>1$ 时，$f(x)>0$。
\end{enumerate}

\item （12分）已知函数 $f(x)=\dfrac{1}{3}x^3-ax^2+1$（$a\in\mathbb{R}$）。
\begin{enumerate}[label=(\arabic*)]
\item 若 $f(x)$ 在 $x=1$ 处取得极值，求 $a$ 的值；
\item 在（1）的条件下，求 $f(x)$ 在 $[-1,3]$ 上的最大值和最小值；
\item 若 $f(x)$ 在 $(0,2)$ 上单调递减，求 $a$ 的取值范围。
\end{enumerate}

\item （12分）已知函数 $f(x)=e^x-ax-1$（$a\in\mathbb{R}$）。
\begin{enumerate}[label=(\arabic*)]
\item 当 $a=1$ 时，求 $f(x)$ 的最小值；
\item 当 $a=2$ 时，证明 $f(x)\geq 0$ 对一切 $x\in\mathbb{R}$ 成立；
\item 若 $f(x)$ 有两个零点，求 $a$ 的取值范围。
\end{enumerate}

\item （12分）已知函数 $f(x)=\ln x-ax+\dfrac{a}{x}-1$（$a\in\mathbb{R}$，$x>0$）。
\begin{enumerate}[label=(\arabic*)]
\item 当 $a=1$ 时，求 $f(x)$ 的单调区间；
\item 当 $a=1$ 时，证明 $f(x)\leq 0$；
\item 当 $a>0$ 时，求 $f(x)$ 的极值点。
\end{enumerate}

\item （14分）已知函数 $f(x)=x\ln x$，$g(x)=x^2-ax+1$（$a\in\mathbb{R}$）。
\begin{enumerate}[label=(\arabic*)]
\item 求 $f(x)$ 的单调区间和最小值；
\item 证明：当 $x>0$ 时，$\ln x\geq 1-\dfrac{1}{x}$；
\item 若对任意 $x\in(0,+\infty)$，$f(x)\geq g(x)$ 恒成立，求 $a$ 的取值范围；
\item 在（3）的条件下，证明：$\displaystyle\sum_{k=1}^{n}\dfrac{1}{k}\geq\ln(n+1)$（$n\in\mathbb{N}^*$）。
\end{enumerate}
\end{enumerate}

\end{document}
